% Copyright (c) 2026 Benjamin Stanley Frohman. The paper is licensed under CC BY 4.0. Lean code is Apache-2.0.
\documentclass[11pt,a4paper]{article}
\usepackage[margin=1in]{geometry}
\usepackage{amsmath,amssymb,amsthm}
\usepackage{hyperref}

\newtheorem{theorem}{Theorem}[section]

\title{Conditional Continuation Theorem for 3D Navier--Stokes:\\Beale--Kato--Majda Criterion and Parabolic Bootstrap}
\author{Benjamin Stanley Frohman\\
\small Independent Researcher, Texas, United States\\
\small ORCID \href{https://orcid.org/0009-0006-7068-3718}{0009-0006-7068-3718}}
\date{\today}

\begin{document}
\maketitle
\noindent{\small Copyright (c) 2026 Benjamin Stanley Frohman. The paper (this PDF and its \LaTeX\ source) is licensed under CC BY~4.0. Lean code in the companion repository is licensed under the Apache License~2.0.}\par\medskip

\begin{abstract}
We record the classical Beale--Kato--Majda continuation criterion in self-contained form. For a maximal Kato solution of the 3D incompressible Navier--Stokes equations on \(\mathbb{T}^3\), integrability of the vorticity supremum in time implies that the solution extends globally and is smooth for positive times. (The classical criterion is stated here in the one-directional continuation form that the argument proves.) The argument uses only the logarithmic Sobolev inequality, Kato--Ponce commutators, a log-Gronwall estimate, and standard parabolic bootstrap; it contains no circularity and no auxiliary fields with incorrect signs.
\end{abstract}

\section{Statement}

\begin{theorem}[Beale--Kato--Majda continuation]
Let \(s>5/2\) and let \(u_0\in H^s(\mathbb{T}^3)\) be divergence-free. Let \(u\in C([0,T^*);H^s(\mathbb{T}^3))\) be the unique maximal local solution of the 3D incompressible Navier--Stokes equations on \(\mathbb{T}^3\). If
\begin{equation}\label{eq:formB}
\int_0^{T^*}\|\omega(t)\|_{L^\infty(\mathbb{T}^3)}\,dt<\infty,
\end{equation}
where \(\omega=\nabla\times u\), then \(T^*=+\infty\) and \(u\in C^\infty(\mathbb{T}^3\times(0,\infty))\).
\end{theorem}

\section{Proof}

\subsection{Logarithmic Sobolev inequality}
The velocity gradient is given by a Calder\'on--Zygmund singular integral of the vorticity. The operator is not bounded on \(L^\infty\). The sharp substitute is the logarithmic estimate (Beale--Kato--Majda / Kozono--Taniuchi): for \(s>5/2\),
\begin{equation}\label{eq:log}
\|\nabla u\|_{L^\infty}
\le C_s\bigl(1+\|\omega\|_{L^\infty}(1+\log^+\|\omega\|_{H^s})+\|\omega\|_{L^2}\bigr).
\end{equation}
(The proof proceeds by Littlewood--Paley decomposition: low frequencies are controlled by the energy, medium frequencies produce a factor equal to the number of shells times \(\|\omega\|_{L^\infty}\), and the high-frequency tail is estimated by Sobolev embedding after a cutoff at frequency comparable to \(\|\omega\|_{H^s}\).)

\subsection{\(H^s\) energy estimate}
Apply \(\Lambda^s=(I-\Delta)^{s/2}\) to the vorticity equation and take the \(L^2\) inner product with \(\Lambda^s\omega\). The viscous term is non-negative. The nonlinear terms are controlled by the Kato--Ponce commutator estimate, yielding
\begin{equation}\label{eq:Hs}
\frac{d}{dt}\|\omega\|_{H^s}^2\le C_s\|\nabla u\|_{L^\infty}\|\omega\|_{H^s}^2.
\end{equation}

\subsection{Log-Gronwall inequality}
Substitute \eqref{eq:log} into \eqref{eq:Hs}. After absorbing lower-order terms one obtains
\begin{equation}
\frac{d}{dt}\|\omega\|_{H^s}^2
\le C\bigl(1+\|\omega\|_{L^\infty}\bigr)\|\omega\|_{H^s}^2\bigl(1+\log^+\|\omega\|_{H^s}\bigr).
\end{equation}
Set \(Y(t)=\log(e+\|\omega(t)\|_{H^s}^2)\). Then
\begin{equation}
\frac{d}{dt}Y(t)\le C\bigl(1+\|\omega(t)\|_{L^\infty}\bigr)Y(t).
\end{equation}
Integration and the ordinary Gronwall lemma produce the double-exponential bound
\begin{equation}
\|\omega(T)\|_{H^s}^2
\le\bigl(e+\|\omega_0\|_{H^s}^2\bigr)^{\exp\bigl(C\int_0^T(1+\|\omega\|_{L^\infty})\,dt\bigr)}.
\end{equation}
Under hypothesis \eqref{eq:formB} the integral is finite, so
\[
\sup_{t\in[0,T^*)}\|\omega(t)\|_{H^s}<\infty.
\]

\subsection{Continuation and parabolic bootstrap}
Kato's local existence theorem asserts that a state with \(\|u\|_{H^s}\le K\) may be continued for a positive time of order \(1/K\). Because the \(H^s\) norm remains bounded up to \(T^*\), the solution can be extended past any finite \(T^*\), a contradiction unless \(T^*=+\infty\).

Once the solution is global and bounded in \(H^s\), Duhamel's formula for the Stokes semigroup
\[
u(t)=e^{\nu t\Delta}u_0-\int_0^t e^{\nu(t-\tau)\Delta}\mathbb{P}\bigl[(u\cdot\nabla)u\bigr](\tau)\,d\tau
\]
together with the smoothing properties of the heat kernel yields, by induction, that \(u(t)\in H^{s+k}\) for every \(k\) and every \(t>0\). Sobolev embedding then gives \(u\in C^\infty(\mathbb{T}^3\times(0,\infty))\).

\end{document}
